\clearpage\fancyhead[L]{\small COMET.}\hypertarget{native-408-part-lower}{}\label{native-408-part-lower}
\begin{center}\large COMET.\end{center}
\begin{center}\hrule\vspace{4pt}\fontsize{10}{12}\selectfont Göttingische gelehrte Anzeigen. No. 60. pp. 593..594. 13 April 1818.\par\vspace{4pt}\hrule\vspace{4pt}\end{center}
\par\indent To make it easier to observe the comet discovered by Mr. \textsc{Pons} at the end of last year, Councillor \textsc{Gauss} has calculated its path in advance for the rest of the time it will be visible. We believe that publishing these results will be of service to observers. The parabolic elements from which this ephemeris was calculated are based on Dr. \textsc{Olbers}’s observations from 3 to 28 March and are as follows:

\begin{center}\fontsize{10}{12}\selectfont
Time of passage through perihelion
\par 26 Feb. 1818. $21^{\mathrm u}\;34^{\mathrm m}\;42^{\mathrm s}$ time on the meridian of Göttingen
\par\smallskip\begin{tabular}{lrr}
Least distance from the Sun & $\cdot\;\cdot\;\cdot\;\cdot$ & $1.20053$\\
Ascending node & $\cdot\;\cdot\;\cdot\;\cdot\;\cdot\;\cdot$ & $70^\circ\;5'\;12''$\\
Inclination of the orbit & $\cdot\;\cdot\;\cdot\;\cdot\;\cdot\;\cdot\;\cdot$ & $90\;0\;0$\\
Longitude of perihelion & $\cdot\;\cdot\;\cdot\;\cdot\;\cdot\;\cdot$ & $183\;22\;58$\\
\end{tabular}\end{center}

\SourceLeaf{409}\par\indent Only when these elements are refined using further observations will it become clear whether the inclination departs appreciably from a right angle, and in which direction. It remains undetermined for now whether the motion is prograde or retrograde with respect to the ecliptic. The time for the ephemeris is $14^{\mathrm u}\,20^{\mathrm m}$ on the meridian of Göttingen. [The Dphemeris is printed in the preceding paper.]

\par\indent Although the comet is already far past perihelion, its light will continue to increase because its distance from Earth is still decreasing, so that on 12 May it will be twice as bright as on 13 March. Since its light is in itself so extraordinarily faint, however, the comet will be observable only through good telescopes.

\par\indent The observations of this comet made at our observatory will be published on another occasion.
